\chapter{Theoretical Background}
\label{ch:theory}

This chapter develops the mathematics of the two architectures. The treatment is deliberately
explicit, because the implementations in Chapter~\ref{ch:implementation} follow these equations
term by term, and the verification in Chapter~\ref{ch:verification} checks them numerically.

\section{Problem formalisation}

Let $\mathcal{V}$ be a finite vocabulary and let a document be a sequence of tokens
$x = (x_1, x_2, \dots, x_n)$ with $x_t \in \mathcal{V}$. Let $\mathcal{C} = \{1, \dots, K\}$ be a
set of $K$ class labels; here $K = 4$.

The task is to learn a function $f_\theta : \mathcal{V}^{*} \rightarrow \Delta^{K-1}$, mapping a
token sequence of arbitrary length to a probability distribution over the $K$ classes, where
$\Delta^{K-1}$ denotes the probability simplex. Given a labelled training set
$\mathcal{D} = \{(x^{(i)}, y^{(i)})\}_{i=1}^{N}$, the parameters $\theta$ are chosen to minimise
the categorical cross-entropy
\begin{equation}
\mathcal{L}(\theta) = -\frac{1}{N} \sum_{i=1}^{N} \log f_\theta(x^{(i)})_{y^{(i)}} .
\label{eq:xent}
\end{equation}

Both architectures studied here share a common three-stage structure:
\begin{enumerate}
\item an \textbf{embedding} stage mapping tokens to dense vectors,
\item a \textbf{contextualisation} stage producing position-wise representations that depend on the
      rest of the sequence --- this is the stage in which the two architectures differ,
\item a \textbf{pooling and classification} stage reducing the variable-length sequence of
      representations to a single vector and projecting it to $K$ logits.
\end{enumerate}

Holding stages 1 and 3 as similar as possible while varying stage 2 is precisely the experimental
design described in Chapter~\ref{ch:method}.

\section{Embedding}

An embedding matrix $E \in \mathbb{R}^{|\mathcal{V}| \times d}$ maps each token index to a row
vector. For a sequence $x$ of length $n$ this yields
\begin{equation}
X = \big(E_{x_1}, E_{x_2}, \dots, E_{x_n}\big)^{\!\top} \in \mathbb{R}^{n \times d}.
\end{equation}
In this study $E$ is randomly initialised and learned jointly with the rest of the model; no
pretrained vectors are used. A designated padding index is held at the zero vector and excluded
from gradient updates.

\section{Recurrent networks and the vanishing gradient problem}
\label{sec:vanishing}

\subsection{The simple recurrent formulation}

A simple recurrent network maintains a hidden state updated at each timestep:
\begin{equation}
h_t = \phi\big(W x_t + U h_{t-1} + b\big), \qquad h_0 = \mathbf{0},
\label{eq:rnn}
\end{equation}
where $\phi$ is a saturating nonlinearity, typically $\tanh$.

\subsection{Backpropagation through time}

Training requires the gradient of a loss $\mathcal{L}$ defined at the final timestep with respect
to parameters that are shared across all timesteps. Applying the chain rule across the unrolled
computation gives
\begin{equation}
\frac{\partial \mathcal{L}}{\partial h_k}
  = \frac{\partial \mathcal{L}}{\partial h_n}
    \prod_{t=k+1}^{n} \frac{\partial h_t}{\partial h_{t-1}} .
\label{eq:bptt}
\end{equation}
From Equation~\ref{eq:rnn}, each Jacobian factor is
\begin{equation}
\frac{\partial h_t}{\partial h_{t-1}}
  = U^{\!\top} \operatorname{diag}\!\big(\phi'(a_t)\big),
  \qquad a_t = W x_t + U h_{t-1} + b .
\end{equation}

\subsection{The pathology}

Taking norms in Equation~\ref{eq:bptt} and writing $\gamma = \sup_z |\phi'(z)|$ --- which equals
$1$ for $\tanh$ and $\tfrac{1}{4}$ for the logistic sigmoid --- gives the bound
\begin{equation}
\left\| \frac{\partial h_t}{\partial h_{t-1}} \right\| \;\le\; \gamma \, \|U\| ,
\qquad\text{hence}\qquad
\left\| \prod_{t=k+1}^{n} \frac{\partial h_t}{\partial h_{t-1}} \right\|
  \;\le\; \big(\gamma \|U\|\big)^{\,n-k} .
\label{eq:vanish}
\end{equation}

The product is geometric in the distance $n-k$. If $\gamma\|U\| < 1$ the gradient decays
exponentially and the network cannot assign credit across long spans --- the \emph{vanishing
gradient} problem. If $\gamma\|U\| > 1$ it grows exponentially --- the \emph{exploding gradient}
problem, which is the more tractable of the two and is routinely handled by gradient-norm
clipping, as is done in this study.

Vanishing gradients cannot be addressed by clipping, because the information is destroyed rather
than merely rescaled. The resolution requires an architectural change.

\section{Long Short-Term Memory}

\subsection{Gating equations}

LSTM introduces a cell state $c_t$ alongside the hidden state $h_t$, together with three gates. At
each timestep:
\begin{align}
i_t &= \sigma\!\left(W_i x_t + U_i h_{t-1} + b_i\right) && \text{input gate} \label{eq:lstm_i}\\
f_t &= \sigma\!\left(W_f x_t + U_f h_{t-1} + b_f\right) && \text{forget gate} \label{eq:lstm_f}\\
g_t &= \tanh\!\left(W_g x_t + U_g h_{t-1} + b_g\right) && \text{candidate cell} \label{eq:lstm_g}\\
o_t &= \sigma\!\left(W_o x_t + U_o h_{t-1} + b_o\right) && \text{output gate} \label{eq:lstm_o}\\
c_t &= f_t \odot c_{t-1} + i_t \odot g_t && \text{cell update} \label{eq:lstm_c}\\
h_t &= o_t \odot \tanh(c_t) && \text{hidden output} \label{eq:lstm_h}
\end{align}
where $\sigma$ is the logistic sigmoid and $\odot$ denotes elementwise multiplication.

Interpretively: the forget gate $f_t$ decides what proportion of the existing cell contents to
retain; the input gate $i_t$ decides how much of the new candidate $g_t$ to admit; and the output
gate $o_t$ decides how much of the resulting cell state to expose as the hidden output.

\subsection{Why this solves the gradient problem}

Consider the derivative of the cell state with respect to its own previous value. Differentiating
Equation~\ref{eq:lstm_c}, and noting that $f_t$, $i_t$ and $g_t$ all depend on $h_{t-1}$ rather
than on $c_{t-1}$ directly, the dominant term is
\begin{equation}
\frac{\partial c_t}{\partial c_{t-1}} \;=\; \operatorname{diag}(f_t) \;+\; (\text{terms via } h_{t-1}).
\label{eq:lstm_grad}
\end{equation}

The leading term is multiplication by the forget gate, with \emph{no weight matrix and no
saturating nonlinearity in the path}. When $f_t \approx \mathbf{1}$, the gradient flows backwards
across that timestep essentially unchanged. Over $n-k$ steps the decay is governed by
$\prod_t f_t$ rather than by $(\gamma\|U\|)^{n-k}$, and since the forget gate is learned, the
network can \emph{choose} to keep this path open for exactly those features requiring long-range
credit assignment. This is the constant error carousel of the original formulation
\cite{hochreiter1997}, and it is the entire reason LSTM works.

\subsection{Fused gate computation}

Equations~\ref{eq:lstm_i}--\ref{eq:lstm_o} share the same two inputs, $x_t$ and $h_{t-1}$. In
implementation the four gate pre-activations are therefore computed in a single pair of matrix
products,
\begin{equation}
\begin{bmatrix} a_i \\ a_f \\ a_g \\ a_o \end{bmatrix}
  = \underbrace{W x_t}_{\in \mathbb{R}^{4d_h}} + \underbrace{U h_{t-1}}_{\in \mathbb{R}^{4d_h}} + b,
\end{equation}
with $W \in \mathbb{R}^{4d_h \times d}$ and $U \in \mathbb{R}^{4d_h \times d_h}$, after which the
result is split into four equal blocks. This is a pure efficiency transformation --- it is
mathematically identical to the separate form --- but it is the form used in the implementation,
and the \emph{ordering} of the four blocks is a convention that must match the reference
implementation for the verification in Chapter~\ref{ch:verification} to be meaningful.

\subsection{Parameter count}

An LSTM layer with input dimension $d$ and hidden dimension $d_h$ has
\begin{equation}
P_{\text{LSTM}} = 4\big(d\,d_h + d_h^2 + d_h\big)
\end{equation}
parameters. A bidirectional layer doubles this.

\section{Bidirectional recurrence}

A bidirectional layer runs two independent recurrences,
\begin{equation}
\overrightarrow{h}_t = \text{LSTM}^{\rightarrow}(x_t, \overrightarrow{h}_{t-1}), \qquad
\overleftarrow{h}_t  = \text{LSTM}^{\leftarrow}(x_t, \overleftarrow{h}_{t+1}),
\end{equation}
and concatenates them, $h_t = [\overrightarrow{h}_t ; \overleftarrow{h}_t] \in \mathbb{R}^{2d_h}$.
Each position is thereby contextualised by the whole sequence rather than only its prefix.

\section{Pooling}
\label{sec:pooling}

Classification requires reducing $(h_1,\dots,h_n)$ to a single vector. Three strategies are
standard.

\keypoint{Last state.} Use $h_n$. For a unidirectional recurrent model this is the natural
summary, but it is sensitive to the final token and discards earlier evidence except insofar as the
state retained it.

\keypoint{Mean pooling.} Use $\frac{1}{n}\sum_t h_t$. This weights all positions equally and is
robust, but for a recurrent model it discards the privileged endpoint that the recurrence was
constructed to produce.

\keypoint{Concatenation.} Use $\big[\tfrac{1}{n}\sum_t h_t \,;\, \overrightarrow{h}_n \,;\,
\overleftarrow{h}_1\big]$, combining both. This is the strategy adopted here, on the reasoning that
the two summaries carry complementary information and documents of 39 tokens are short enough that
the endpoint alone is noisy.

With padded batches, all pooling operations must be masked --- a point of considerable practical
importance developed in Section~\ref{sec:masking}.

\section{Attention}

\subsection{Scaled dot-product attention}

Given queries $Q \in \mathbb{R}^{n \times d_k}$, keys $K \in \mathbb{R}^{m \times d_k}$ and values
$V \in \mathbb{R}^{m \times d_v}$, scaled dot-product attention is
\begin{equation}
\operatorname{Attention}(Q,K,V)
  = \operatorname{softmax}\!\left(\frac{QK^{\top}}{\sqrt{d_k}}\right) V .
\label{eq:attention}
\end{equation}

The matrix $QK^{\top} \in \mathbb{R}^{n\times m}$ contains a compatibility score for every
query--key pair. The row-wise softmax converts each row into a probability distribution over the
$m$ positions, and the product with $V$ forms, for each query, a convex combination of value
vectors.

In \emph{self}-attention, $Q$, $K$ and $V$ are all linear projections of the same input sequence,
so $n = m$ and every position attends to every position including itself.

\subsection{Why the $\sqrt{d_k}$ scaling is necessary}

The scaling factor is not cosmetic. Suppose the components of $q$ and $k$ are independent with zero
mean and unit variance. Then
\begin{equation}
\operatorname{E}\!\left[q \cdot k\right] = 0,
\qquad
\operatorname{Var}\!\left[q \cdot k\right]
  = \operatorname{Var}\!\left[\sum_{j=1}^{d_k} q_j k_j\right] = d_k .
\end{equation}

The dot products therefore have standard deviation $\sqrt{d_k}$, growing with the head dimension.
For large $d_k$ the softmax receives inputs of large magnitude and saturates: its output approaches
a one-hot vector, and its Jacobian approaches zero. Writing $p = \operatorname{softmax}(s)$, the
Jacobian is
\begin{equation}
\frac{\partial p_i}{\partial s_j} = p_i\big(\delta_{ij} - p_j\big),
\end{equation}
which vanishes as any $p_i \to 1$. Dividing by $\sqrt{d_k}$ restores unit variance to the scores
and keeps the softmax in its responsive region, preserving gradient flow.

\subsection{Multi-head attention}

A single attention operation produces one weighted average per position, which constrains the model
to one relational pattern per layer. Multi-head attention runs $h$ attention operations in
parallel on different learned projections:
\begin{align}
\operatorname{head}_j &= \operatorname{Attention}\!\left(XW_j^{Q},\, XW_j^{K},\, XW_j^{V}\right), \\
\operatorname{MHA}(X) &= \left[\operatorname{head}_1 ; \dots ; \operatorname{head}_h\right] W^{O} .
\label{eq:mha}
\end{align}
With $d_k = d_v = d_{\text{model}}/h$, total computation is approximately that of a single head at
full dimension, but the model can attend to several distinct relational patterns simultaneously.

In implementation the per-head projections are not separate matrices. A single projection of width
$d_{\text{model}}$ is computed and then \emph{reshaped} into $h$ blocks of width $d_k$. This
reshape is the single most error-prone step in implementing attention, and it is the specific
target of verification check~3 in Chapter~\ref{ch:verification}.

\subsection{Masking}
\label{sec:masking}

Batched processing requires padding sequences to a common length. Padded positions must not
influence the result. In attention this is enforced by adding a mask to the scores before the
softmax:
\begin{equation}
s'_{ij} = \begin{cases}
s_{ij} & \text{if position } j \text{ is a real token} \\
-\infty & \text{if position } j \text{ is padding,}
\end{cases}
\end{equation}
so that $\operatorname{softmax}(s')$ assigns exactly zero weight to padded keys.

In practice $-\infty$ is replaced by a large finite negative constant. The reason is numerical: if
an entire row were masked, $\operatorname{softmax}$ of a vector of $-\infty$ evaluates to
$0/0 = \text{NaN}$, which propagates through the network and destroys the gradients. A large
finite constant yields a harmless uniform distribution in that degenerate case.

\section{Positional encoding}

Self-attention as defined in Equation~\ref{eq:attention} is permutation-equivariant: permuting the
input rows permutes the output rows identically, with no other change. A model built from
self-attention alone therefore cannot distinguish word order at all --- it would be a bag-of-words
model with a sophisticated weighting scheme.

Position must consequently be injected explicitly. The sinusoidal scheme of
Vaswani et al.\ \cite{vaswani2017} defines
\begin{align}
PE_{(\text{pos},\,2i)}   &= \sin\!\left(\frac{\text{pos}}{10000^{2i/d_{\text{model}}}}\right), \\
PE_{(\text{pos},\,2i+1)} &= \cos\!\left(\frac{\text{pos}}{10000^{2i/d_{\text{model}}}}\right),
\end{align}
and adds $PE$ to the token embeddings. Each dimension is a sinusoid of a different wavelength,
forming a geometric progression from $2\pi$ to $10000\cdot 2\pi$. The scheme has two useful
properties: it requires no parameters, and because $PE_{\text{pos}+k}$ is a fixed linear function
of $PE_{\text{pos}}$ for any fixed offset $k$, relative position is linearly recoverable.

Token embeddings are scaled by $\sqrt{d_{\text{model}}}$ before the positional signal is added, so
that the two contributions are of comparable magnitude at initialisation.

\section{Layer normalisation}

Layer normalisation \cite{ba2016} standardises each position's feature vector independently:
\begin{equation}
\operatorname{LN}(x) = \gamma \odot \frac{x - \mu}{\sqrt{\varsigma^2 + \epsilon}} + \beta,
\qquad
\mu = \frac{1}{d}\sum_{j=1}^{d} x_j,
\qquad
\varsigma^2 = \frac{1}{d}\sum_{j=1}^{d} (x_j - \mu)^2 ,
\end{equation}
with learned $\gamma, \beta \in \mathbb{R}^{d}$ and a small $\epsilon$ for numerical stability.

Unlike batch normalisation, the statistics are computed over the feature dimension of a single
example, so the operation is independent of batch composition and of sequence length. Note that
$\varsigma^2$ is the \emph{biased} variance, normalising by $d$ rather than $d-1$; this detail
matters for the numerical agreement verified in Chapter~\ref{ch:verification}.

\section{Residual connections and block ordering}

Each sublayer is wrapped in a residual connection, which provides an identity path for gradients.
Two orderings are in use:
\begin{align}
\text{post-norm:}\quad & x \leftarrow \operatorname{LN}\!\big(x + \operatorname{Sublayer}(x)\big), \\
\text{pre-norm:}\quad  & x \leftarrow x + \operatorname{Sublayer}\!\big(\operatorname{LN}(x)\big).
\end{align}

In the post-norm arrangement the residual path passes through a normalisation at every layer, and
at initialisation the expected gradient magnitude scales unfavourably with depth, which in practice
necessitates a learning-rate warmup schedule. In the pre-norm arrangement the residual path is
unobstructed from input to output, gradients are better conditioned at initialisation, and warmup
becomes unnecessary \cite{xiong2020}. The present study uses pre-norm, since the training budget
is a small number of epochs and a warmup schedule would consume a meaningful fraction of it.

The complete encoder block is
\begin{align}
x &\leftarrow x + \operatorname{Dropout}\big(\operatorname{MHA}(\operatorname{LN}(x))\big), \\
x &\leftarrow x + \operatorname{Dropout}\big(W_2\,\operatorname{GELU}(W_1 \operatorname{LN}(x))\big).
\end{align}

\section{Complexity analysis}

Table~\ref{tab:complexity} summarises the structural comparison, following the analysis of
Vaswani et al.\ \cite{vaswani2017}. Let $n$ be sequence length and $d$ the representation
dimension.

\begin{table}[ht]
\centering
\caption[Complexity comparison of recurrent and self-attentive layers]{Per-layer complexity,
sequential depth and maximum path length. \emph{Sequential operations} counts steps that cannot be
parallelised; \emph{maximum path length} is the number of computational steps separating the two
most distant positions.}
\label{tab:complexity}
\begin{tabular}{lccc}
\toprule
\textbf{Layer type} & \textbf{Complexity per layer} & \textbf{Sequential operations} & \textbf{Max path length} \\
\midrule
Recurrent (LSTM)    & $O(n \cdot d^2)$ & $O(n)$ & $O(n)$ \\
Self-attention      & $O(n^2 \cdot d)$ & $O(1)$ & $O(1)$ \\
\bottomrule
\end{tabular}
\end{table}

Three consequences follow, and all three are tested empirically in Chapter~\ref{ch:results}.

\keypoint{Path length.} Self-attention connects any two positions in one step; recurrence requires
$O(n)$ steps. This is the basis of the Transformer's advantage on long-range dependency. Its
practical value is proportional to how much long-range structure the task actually contains.

\keypoint{Parallelism.} Recurrence requires $O(n)$ strictly sequential steps; self-attention
requires $O(1)$. On parallel hardware this is a decisive advantage in wall-clock terms.

\keypoint{Arithmetic.} The comparison between $O(n d^2)$ and $O(n^2 d)$ depends on the ratio
$n/d$. Self-attention performs \emph{less} total arithmetic only when $n < d$. With $n = 64$ and
$d = 128$, as in this study, the two are of the same order, and neither dominates. This observation
is central to the cost result reported in Section~\ref{sec:cost}, and it is why that result does
not contradict the theory.
